% ======================================================================
% Journal of Computer-Aided Molecular Design (2026)
% Short Communication - V6.2 (expanded to accepted-paper rigor level)
% ======================================================================
\documentclass[global,twocolumn]{svjour}

\usepackage[T1]{fontenc}
\usepackage{lmodern}
\usepackage{amsmath}
\usepackage{graphicx}
\usepackage{booktabs}
\usepackage{array}
\usepackage[expansion=false]{microtype}
\usepackage{xurl}
\usepackage[numbers,sort&compress]{natbib}
\usepackage[colorlinks=true,linkcolor=blue,citecolor=blue,urlcolor=blue]{hyperref}
\usepackage{caption}
\usepackage{placeins}

\journalname{Journal of Computer-Aided Molecular Design}
\newcommand{\st}{\footnotesize\setlength{\tabcolsep}{3pt}}

\begin{document}

\title{Cognate redocking as a quality-control step for cryo-EM kinase
       receptor structures: a study of ATR/CHK1 and
       CDK-activating kinase}
\titlerunning{Cognate redocking for cryo-EM kinase structures}

\author{Ethan Ye \and Aryan Padarthi \and Sungju Kim \and Alexandre Elie-Dit-Cosaque \and Eunho Jo}
\authorrunning{Ye et al.}
\institute{Allen High School, Allen ISD, Allen, TX 75013, USA \\
           \email{aryan.padarthi@student.allenisd.org}}
\offprints{Aryan Padarthi}

\date{Received: / Accepted:}
\maketitle

% ======================================================================
\begin{abstract}
Reproducing a receptor's own deposited experimental ligand model (cognate
redocking) is a quality check commonly applied in structure-based virtual
screening. The cryo-EM ATR structures bound to berzosertib (PDB~9L40,
2.87~\AA{}) and camonsertib (PDB~9L4B, 3.20~\AA{}) were redocked with
AutoDock Vina (2.0~\AA{} symmetry-corrected RMSD pass threshold): 9L4B
passed (0.61$\pm$0.02~\AA{}) while the higher-resolution 9L40 failed
(2.53$\pm$0.01~\AA{}), not explained by the nearby allosteric pocket
(60.65~\AA{} away). The procedure was repeated on 14 cryo-EM structures of
the CDK7/cyclin~H/MAT1 CAK complex (a single-target series): again, 13/14
structures did not pass. This is not a comprehensive benchmark, but it does
serve as a pipeline sanity check: the same pipeline was applied to 20
objectively chosen CASF-2016 complexes (8/20 passing, 40\%; exact 95\%~CI
19--64\%), arguing against a badly broken pipeline. Neither nominal
resolution nor the wwPDB map/model metrics (Q-score, residue inclusion)
correlated with redocking performance; given only two Q-scored passing
structures, this is a descriptive finding, not a conclusive statistical
correlation. Against a matched rigid control (4.50$\pm$0.05~\AA{}),
allowing five active-site residues within 4.5~\AA{} of native berzosertib
to move improved the pose of 9L40 (3.00$\pm$0.92~\AA{}) but did not let it
pass. Docking a 40-compound pilot library against excluded 9L40 instead of
validated 9L4B left the top five compounds unchanged but reshuffled the
screen's mid-ranks (Spearman $\rho=0.21$; only 5/10 shared top-10
compounds). Cognate redocking is a structure-specific quality-control step
whose necessity and sufficiency for prospective enrichment require
systematic, multi-target validation.
\end{abstract}

\keywords{molecular docking \hspace{0.2em} redocking validation \hspace{0.2em} cryo-EM \hspace{0.2em} ATR kinase \hspace{0.2em} virtual screening \hspace{0.2em} AutoDock Vina \hspace{0.2em} BRICS \hspace{0.2em} CASF-2016}

% ======================================================================
\section{Introduction}

Molecular docking is one of the most commonly employed computational tools
in structure-based drug discovery because it provides a good compromise
between speed and accuracy~\cite{trott2010,hartshorn2007}. Several factors
are crucial for the reliability of the results obtained: receptor
structure, ligand preparation, scoring function, search algorithm, and
validation protocol. One of the most critical among these is the quality
and conformation of the receptor structure, since all downstream analyses
will inherit any errors introduced at this stage. This article aims to
test that claim directly.

Cryo-EM has become an effective method for determining macromolecular
structures, competing in many cases with X-ray crystallography at
increasingly high resolution~\cite{pintilie2020,kuhlbrandt2014}. Although
resolution and a deposited structure's suitability for computational
docking might be expected to track together, recent work indicates that
medium-resolution cryo-EM structures, cross-docked under a standard
protocol, succeed at only about half the rate of high-resolution crystal
structures~\cite{lee2023}. Several recent studies have also emphasized the
importance of per-structure validation using AlphaFold2-predicted
models~\cite{holcomb2023,scardino2022}: structure-specific empirical
testing may be required regardless of how high-quality the input model
appears, or its source.

ATR and CHK1 are ATP-competitive kinases with several chemical inhibitors
reported for structure-based
design~\cite{kumagai2006,bass2016,saldivar2017,casper2002,debatisse2012,fokas2012}.
Recently published cryo-EM structures of the ATR kinase domain (PDB~9L40,
bound to berzosertib, 2.87~\AA{}; PDB~9L4B, bound to camonsertib,
3.20~\AA{})~\cite{wang2025} provide better starting points for docking than
the original 4.7~\AA{} ATR--ATRIP complex (PDB~5YZ0). In a companion study,
berzosertib was shown to engage an additional dimer-interface allosteric
site~\cite{wang2025}, further highlighting the intrinsic complexity of this
system. The challenge posed by these data is practical, at the level of
experimental protocol, rather than target biology: docking to a deposited
cryo-EM coordinate set treats those coordinates as the truth, but if that
single snapshot model does not reproduce the ligand pose on which it was
based, no amount of scoring innovation can restore a meaningful answer. We
therefore asked whether nominal resolution and standard cryo-EM map-model
quality metrics are sufficient to determine whether these deposited
structures can reproduce their own experimental ligand poses under a fixed
docking protocol.

Reproducing a receptor's own bound ligand pose (cognate redocking) is an
important sanity check for any docking methodology, but it does not
necessarily indicate predictive capability for prospective
screening~\cite{jain2008}. Our work should be read as a methodological
investigation into this process, not a broad virtual screening exercise. We
also remind the reader that predicted binding, while useful, cannot provide
evidence of functional modulation~\cite{kumagai2006,bass2016}.

% ======================================================================
\section{Materials and methods}
\label{sec:methods}

\subsection{Receptor structures and preparation}

Three receptor structures were retrieved in mmCIF format from the RCSB PDB
AWS Open Data mirror (snapshot~2026-01-01): PDB~9L40 (ATR, bound to
berzosertib, 2.87~\AA{}), PDB~9L4B (ATR, bound to camonsertib, 3.20~\AA{}),
and PDB~2YM8 (human CHK1, 2.07~\AA{})~\cite{wang2025}. We include 2YM8 as a
non-cryo-EM (X-ray crystallographic) kinase reference structure to test
whether our findings are reproducible on non-cryo-EM-derived receptors;
with only one such structure to compare against, this can only be treated
as a hypothesis ($n=1$), which we report as such (Limitations). PDB~5YZ0
(4.7~\AA{}) was omitted due to its poor resolution and consequent inability
to reliably resolve side chains within the ATP pocket~\cite{wang2025}. Both
protein chains (full homodimer) were retained using Biopython~\cite{cock2009}.

For 9L40, contact analysis (heavy atoms within 4.5~\AA{} of each protein
chain) identified residues A2701/B2701 as the ATP-competitive active site
and A2702/B2702 as the dimer-interface allosteric site~\cite{wang2025}. The
centroid-to-centroid distance between these sites is 60.65~\AA{}, almost
three times the 22~\AA{} docking box used here, so the ATP-pocket redocking
result reported below cannot have been affected by the allosteric site.

\textbf{Receptor preparation} was performed with Open Babel~3.1.1~\cite{oboyle2011}
and Meeko (2025). For each structure:
\begin{itemize}\sloppy
\item Hydrogens were added at pH~7.4 (Open Babel \texttt{--pH 7.4});
\item Alternate locations were resolved by keeping the highest-occupancy
  conformer; none of the alternate conformations are located near the
  cognate ligand in any of the three primary structures (within 5~\AA{});
\item Crystallographic waters and non-ligand heteroatoms were removed;
\item Asp/Glu negative; Lys/Arg positive; histidine tautomer states were
  determined by Meeko during the flexible-receptor preparation step, which
  requires an explicit HIE assignment for HIS2477 (Section~2.8);
\item Meeko's receptor-preparation step deleted \textasciitilde400~residues
  whose side-chain density could not be matched against a template at
  approximately 3~\AA{} resolution; automatically dropping surface/loop
  side chains that cannot be reliably placed in an incomplete cryo-EM map
  is a standard feature of receptor preparation at this resolution regime,
  not a study-specific deviation, and all deleted residues were $>$10~\AA{}
  from the docking box (Supplementary Table~S2);
\item Proteins were converted to PDBQT with Open Babel using Gasteiger-Marsili
  partial charges~\cite{gas_marsili} (\texttt{--gen3d --partialcharge
  gasteiger}).
\end{itemize}

All receptors and ligands in this study (redocking set, pilot library, CAK
survey, and CASF-2016 subset alike) were docked at physiological pH~7.4 and
using Gasteiger-Marsili charges. These are parameter-free, per-atom default
charge models that do not need any target- or ligand-specific tweaking --
the appropriate choice for a protocol that is deliberately held constant
across a structurally and chemically diverse set of receptors and compounds
rather than tailored specifically for any of them. A more subtle charge
model or per-compound prediction of protonation state would reintroduce
precisely the type of ligand-specific tweaking that this study's design
aims to avoid.

Because all ligands docked in this study are single-site, ATP-competitive
kinase inhibitors of approximately 550~Da or less, the 22~\AA{} cubic
search box, centered on the cognate ligand centroid, was retained as the
sole, unchanged protocol parameter for all receptors and all 45
library/survey ligands (5 seeds, 35 derivatives, and the CASF/CAK cognate
ligands) instead of being sized per ligand: it is large enough to contain
both the ligand itself and its first-shell pocket residues in every
structure tested, which renders a static box size neutral with respect to
the comparison between receptors. This follows the advice of Feinstein and
colleagues~\cite{feinstein2021} that search-box size be rationalized
according to the specific ligand and receptor rather than chosen by
default; no systematic per-ligand box-size sensitivity sweep across the
full benchmark was undertaken, which is noted as a limitation
(Section~3.7).

wwPDB validation reports give per-ligand information for the active-site
ligand copy used as the redocking reference in each structure: 9L40/A2701,
Q-score~0.553, residue inclusion~0.818; 9L4B/A2701, Q-score~0.528, residue
inclusion~0.967~\cite{pintilie2020}. These values indicate that both
deposited ligand models show similar local map support, although a similar
Q-score does not guarantee that the pose itself is
correct~\cite{pintilie2020}.

\subsection{Seed library and derivatives}

Five compounds comprised the seed library: the three deposited active-site
ligands (berzosertib, camonsertib, and the CHK1 inhibitor from 2YM8), plus
prexasertib and ceralasertib (SMILES checked against reported molecular
formula and exact mass). These were split into 22 unique fragments using
the BRICS algorithm~\cite{degen2008}, recombined using RDKit~\cite{rdkit},
and capped at 5,000 candidates. Candidates were pre-filtered by relaxed
Lipinski~\cite{lipinski1997} and Veber~\cite{veber2002} rules, plus
PAINS~\cite{baell2010}/Brenk structural-alert removal~\cite{brenk2008},
leaving 786 candidates. Retained compounds were assessed against four
drug-likeness rules (Lipinski, Veber, Ghose~\cite{ghose1999},
Egan~\cite{egan2000}), of which 680 of 786 (86.5\%) satisfied $\ge$3/4
rules. Derivatives selected for pilot docking ($n=35$) were those
satisfying the most rules, ranked by ascending molecular weight. The
complete ranking and selection scores are in Supplementary Table~S4. This
selection protocol is necessarily biased toward smaller, rule-compliant
molecules, which should be kept in mind when interpreting the pilot docking
ranking.

\subsection{Redocking validation protocol}

Each receptor was validated by extracting its deposited ligand, redocking
into the same receptor/box using AutoDock Vina~1.2.7 (Python
bindings)~\cite{trott2010} exhaustively (exhaustiveness~32), seeded with
three independent replicates (seeds~1000, 1001, 1002), and calculating the
RMSD (heavy atoms) between the redocked pose and the deposited native pose,
using a pre-specified pass threshold of $\le$2.0~\AA{}~\cite{hartshorn2007}.
The threshold was set to conform to the Astex Diverse Set
convention~\cite{hartshorn2007} and was pre-specified prior to testing any
receptor, to prevent post-hoc threshold selection.

Two RMSD measures are reported: naive index-matched RMSD and
symmetry-corrected RMSD (spyrmsd~v0.6.2, graph-isomorphism atom
matching~\cite{meli2020}). The symmetry-corrected RMSD is defined as the
root-mean-square deviation between a docked pose $P_d$ and a reference pose
$P_r$:
\[
\mathrm{RMSD_s}(P_d, P_r) = \min_{\sigma \in S} \sqrt{\frac{1}{N}
\sum_{i=1}^{N} \| \mathbf{x}_{d,i} - \mathbf{x}_{r,\sigma(i)} \|^2 }
\]
where $\sigma$ is a permutation of equivalent atoms in the set $S$ of
topological symmetries and $N$ is the number of heavy atoms. This
correction is useful for compounds like berzosertib that contain internal
topological symmetries (equivalent ring atoms, sulfone oxygens), for which
the naive index-matched RMSD may overstate the true error of the
pose~\cite{meli2020}.

Search reproducibility $\rho_{\text{repro}}$ is the average pairwise RMSD
between the best-scoring poses of each replicate across the three seeded
runs:
\[
\rho_{\text{repro}} = \frac{1}{3} \sum_{i < j} \mathrm{RMSD}(P_{\text{best},i},
P_{\text{best},j})
\]
This metric assesses the random variation in the Vina search rather than
its accuracy against the native pose, a distinction kept separate
throughout this study (Results).

Both sanity controls yielded the expected result, with crystal-vs-itself
giving exactly 0.000~\AA{} and the exact 1.0~\AA{} rigid translation
recovered as exactly 1.000~\AA{}. If a receptor failed the redocking
control in any replicate, it was excluded from candidate scoring.

\subsection{Ligand preparation}

For cognate redocking, the deposited ligand chemical state (connectivity,
stereochemistry, protonation, tautomer) was preserved unchanged, since the
purpose of that test is to ask whether the docking protocol can recover a
known experimental pose, not to explore alternative protonation states. For
the pilot screen, all seed and derivative ligands were prepared by 3D
embedding, protonating at pH~7.4 with Open Babel, assigning Gasteiger-Marsili
partial charges, and converting to PDBQT -- the same parameter-free
treatment used for every receptor and ligand in this study (Section~2.1). A
more detailed per-compound search of tautomers/protomers was not attempted
and is not claimed here.

\subsection{Pilot library docking}

The 5 seeds plus the 35 selected derivatives were docked against each
passing receptor using the identical protocol except at exhaustiveness~16
(three seeds: 2000, 2001, 2002). For consistency of the protocol argument,
we validated redocking at this screening-level exhaustiveness as well: the
same pass/fail pattern held (9L40 3.01$\pm$0.06~\AA{} fail, 9L4B
0.59$\pm$0.02~\AA{} pass, 2YM8 0.92$\pm$0.01~\AA{} pass), indicating that
the reliability conclusion is robust to this choice of parameter.

\subsection{Multi-structure survey}

The identical validation protocol was also applied to 14 additional
cryo-EM structures of the CDK7/cyclin~H/MAT1 CAK complex~\cite{cushing2024},
each bound to a different ATP-competitive inhibitor. All 14 structures come
from one target system and one research group, making this a clustered
single-target set rather than 14 independent targets. Given this close
relationship, the 13/14 result for this set should be viewed as descriptive
evidence of failure within this series, not a population estimate. The
pipeline was verified on 9L4B: it reproduced the original result almost
exactly (RMSD~0.60$\pm$0.03 vs.\ 0.61$\pm$0.02~\AA{}; Q-score~0.528,
identical), confirming that the automated survey pipeline is consistent
with the hand-curated protocol.

\subsection{Positive-control sanity check}

Twenty CASF-2016~\cite{su2019} complexes were selected by an objective,
non-random rule fixed before any of the 20 was inspected individually: the
first 20 PDB IDs, in file order, of a complete CASF-2016 core-set listing
redistributed, not curated or re-ordered by us, in the \texttt{DeepDock}
GitHub repository.\footnote{\url{https://github.com/OptiMaL-PSE-Lab/DeepDock},
file \texttt{Validation\_Docking/DockingResults\_CASF2016\_CoreSet.csv}.}
This sequential rule, rather than a stratified or pseudorandom draw, was
chosen specifically to avoid any appearance of hand-picking targets
favorable to the pipeline. The 285-complex CASF-2016 core
set~\cite{su2019} is a popular test case for docking protocol validation,
and the reported 58\% Vina self-docking success rate on this
benchmark~\cite{zhang2023} is a helpful contextual comparison. That value
was obtained under particular preparation and sampling conditions distinct
from ours, so we do not take it as representative of any ``Vina success
rate.'' This test is meant purely as a check on our pipeline, not as a
broad benchmark of docking accuracy: 20 complexes, under one fixed generic
protocol, cannot replace a systematic evaluation of accuracy, and failing
to reject a null hypothesis of equality to the 58\% literature value does
not itself indicate equality. We report exact binomial 95\% confidence
intervals using the Clopper-Pearson method~\cite{clopper1934}.

\subsection{Flexible-sidechain redocking}

As a test for any potential role of rigid conformations in the failure of
9L40, five active-site residues within 4.5~\AA{} of native berzosertib
(chain~A: LYS2327, TYR2365, GLU2378, TRP2379, HIS2477) were selected
\emph{a priori} (not based on inspection of the redocked poses) and defined
as flexible side chains for Vina. The 4.5~\AA{} distance cutoff matches the
same heavy-atom contact-shell radius used earlier in this study to
differentiate the ATP-competitive site from the allosteric site
(Section~2.1); the five identified residues are therefore a natural,
predetermined consequence of applying this fixed radius to the native pose,
not a deliberate choice for effect.

The flexible and rigid parts of the receptor were prepared with Meeko
(\texttt{mk\_prepare\_receptor.py}), which requires an explicit tautomer
assignment for every histidine, including HIS2477 (HIE). As in
Section~2.1, template-mismatched surface residues were automatically
deleted by Meeko; none of the five targeted flexible residues was among
them. To give the flexible side chains room to explore alternative rotamers
without triggering a ligand-outside-box error, the docking box size was
increased from 22~\AA{} (rigid redocking, Section~2.1) to 30~\AA{}. Because
this changes three things relative to the original rigid result at once --
flexibility, box size, and the Meeko-prepared receptor -- a matched control
was run alongside the flexible run: fully rigid redocking (no residues
declared flexible) with the identical enlarged 30~\AA{} box and the
identical Meeko-prepared receptor and ligand, isolating the effect of
flexibility from the other two changes. Both conditions used three seeded
replicates (1000, 1001, 1002, exhaustiveness~32), consistent with the
original rigid-redocking protocol. Individual replicate RMSDs are provided
in Supplementary Table~S5.

\subsection{Rank concordance analysis}

To provide a fair comparison, the 40-compound pilot set was docked against
excluded 9L40 using the same three seeds as for 9L4B (2000, 2001, 2002).
Rank concordance metrics included: Spearman rank correlation $\rho$ with
bootstrap 95\% confidence interval (10,000 resamples), Kendall~$\tau$,
mean absolute rank displacement, and top-10 overlap fraction. For $\rho$,
the bootstrap approach estimates uncertainty by resampling compound pairs
with replacement~\cite{efron1993}, accounting for the small set size.

\subsection{Computational environment and reproducibility}

All calculations: Vina~1.2.7 (Python bindings), Open Babel~3.1.1,
RDKit~2026.03.6, Biopython~1.88, Meeko (2025), all CPU only on a standard
workstation (no GPU acceleration). Exact command-line arguments, prepared
PDBQT files, and individual replicate RMSDs are provided in the repository.
Supplementary Table~S6 provides the three seeded RMSD values for each
structure individually. Reported mean~$\pm$~SD describes variation among
the three docking runs but is not an inferential measure of uncertainty
about the receptor's suitability. All software versions, random seeds,
thresholds, and selection criteria were recorded prior to the first
execution of the pipeline and are available for independent verification.

% ======================================================================
\section{Results and discussion}
\label{sec:results}

\subsection{Redocking validation}

Redocking validation results are summarized in Table~1. 9L4B and 2YM8
passed comfortably (symmetry-corrected RMSD~0.61$\pm$0.02~\AA{} and
0.76$\pm$0.01~\AA{}). 9L40, the higher-resolution structure, failed
(RMSD~2.53$\pm$0.01~\AA{}). The nearest allosteric-site ligand copy is
60.65~\AA{} from the active site, nearly three times the 22~\AA{} docking
box, which rules out the possibility that the docking box reached the
allosteric site. Per the pre-specified rule, 9L40 was excluded from the
candidate-scoring process. Figure~1 shows the native and docked poses for
both structures.

\begin{figure}[t]
\centering
\includegraphics[width=\columnwidth]{figure_redocking_overlay.png}
\caption{Native (grey) vs. redocked (colored) best pose. Left:
9L40/berzosertib (symmetry-corrected RMSD 2.53~\AA{}). Right:
9L4B/camonsertib (0.61~\AA{}).}
\label{fig:overlay}
\end{figure}

\begin{table}[t]
\centering
\caption{Redocking validation (three seeded Vina replicates,
exhaustiveness~32)}
\label{tab:redocking}
\st
\begin{tabular}{lcccl}
\toprule
Structure & Res. (\AA{}) & RMSD (\AA{}) & Pass? & Repro.\ (\AA{})$^a$ \\
\midrule
9L40 (ATR)      & 2.87 & 2.53$\pm$0.01 & \textbf{No} & 0.56 \\
9L4B (ATR)      & 3.20 & 0.61$\pm$0.02 & Yes         & 0.045 \\
2YM8 (CHK1)$^b$ & 2.07 & 0.76$\pm$0.01 & Yes         & 0.078 \\
\bottomrule
\end{tabular}
\vspace{2pt}
{\footnotesize $^a$ Search reproducibility: mean pairwise RMSD between
replicate best poses. $^b$ Non-cryo-EM (X-ray) reference structure.}
\end{table}
\FloatBarrier

\subsection{Multi-structure survey}

Table~2 reports results for all 17 structures tested. Of these, 14 of 17
(82\%) do not pass the 2.0~\AA{} threshold; only three pass it: 9L4B, 2YM8,
and one CAK entry (8P75). There was no statistically significant monotonic
association between reported resolution and redocking outcome (Spearman
$\rho=-0.28$, $p=0.28$, $n=17$), nor between the evaluated global measure
of residue inclusion and redocking outcome ($r=-0.14$, $p=0.61$, $n=16$).
For the 12 structures with a reported Q-score, the two passing structures
had a lower mean Q-score (0.569) than the 10 failing structures (0.688),
which implies a descriptive trend but is not a definitive statistical
correlation given the small number of Q-scored passing structures
($n=2$); this result should not be taken as an indication that Q-score has
no predictive value.

The CAK failures showed RMSDs in the range 5--7~\AA{} but docked ligand
centroids within 1.1--4.1~\AA{} of their native centroids, suggesting the
docked ligand had been reoriented within the pocket rather than moved far
away from it. This is consistent with orientational degeneracy at the
kinase hinge -- where multiple positions may be reasonable for an
ATP-competitive scaffold -- or with rigid-receptor limitations that prevent
the docking engine from exploring the experimentally observed binding
mode, rather than with a setup problem such as a misplaced box.

The 14 CAK structures come from a single target system and research group;
they are thus clustered observations, not 14 independent targets. Within
this single-target series, 13 of 14 structures failed the prespecified
criterion. Because of this clustering, the 13/14 observation should not be
interpreted as an estimate of the underlying failure probability for
kinase structures generally, but rather as descriptive evidence of failure
within this particular CAK series.

\begin{table}[t]
\centering
\caption{Multi-structure survey, all 17 structures tested, sorted by
symmetry-corrected RMSD. n/a: Q-score not reported for this experimental
method (X-ray) or ligand occupancy state (split-altloc CAK ligands).}
\label{tab:survey}
\st
\begin{tabular}{lccccc}
\toprule
Structure & Res. (\AA{}) & RMSD (\AA{}) & Pass? & Q-score & Res.\ incl. \\
\midrule
9L4B (ATR)  & 3.20 & 0.61 & Yes & 0.528 & 0.967 \\
2YM8 (CHK1) & 2.07 & 0.76 & Yes & n/a   & n/a   \\
8P75 (CAK)  & 2.00 & 1.03 & Yes & 0.609 & 0.724 \\
9L40 (ATR)  & 2.87 & 2.53 & No  & 0.553 & 0.818 \\
8P72 (CAK)  & 1.90 & 2.70 & No  & 0.707 & 0.963 \\
8P78 (CAK)  & 1.90 & 2.98 & No  & n/a   & 0.828 \\
8P6Z (CAK)  & 2.10 & 4.84 & No  & 0.705 & 0.885 \\
8P77 (CAK)  & 1.80 & 5.02 & No  & n/a   & 0.793 \\
8P73 (CAK)  & 2.00 & 5.84 & No  & 0.722 & 0.800 \\
8P6V (CAK)  & 1.90 & 6.14 & No  & 0.738 & 0.966 \\
8P7L (CAK)  & 2.10 & 6.16 & No  & 0.708 & 0.969 \\
8P70 (CAK)  & 2.00 & 6.25 & No  & n/a   & 1.000 \\
8P76 (CAK)  & 2.00 & 6.60 & No  & 0.640 & 0.735 \\
8P6X (CAK)  & 1.90 & 6.60 & No  & n/a   & 0.964 \\
8P74 (CAK)  & 2.20 & 6.67 & No  & 0.628 & 0.793 \\
8P71 (CAK)  & 2.00 & 7.02 & No  & 0.720 & 0.926 \\
8P6W (CAK)  & 1.90 & 7.33 & No  & 0.762 & 0.929 \\
\bottomrule
\end{tabular}
\end{table}
\FloatBarrier

\subsection{Positive-control sanity check}

The unmodified pipeline passed 8 of 20 (40\%) CASF-2016 subset complexes
selected by the first-20-in-file-order rule described in Methods (Table~3).
The exact 95\% binomial confidence interval (Clopper-Pearson) is
19.1--63.9\%. We note that the purpose of this subset is as a pipeline
sanity check, not an exhaustive evaluation of docking accuracy: failure to
reject a null hypothesis does not suggest equivalence, only that we are
unable to detect a significant difference. It indicates that our pipeline
is not badly broken, but not that our pipeline is definitively valid.
Although the observed 40\% is not significantly different from the
literature Vina value of 58\% on the full 285-complex core
set~\cite{zhang2023} (one-sided binomial $p=0.08$), the comparison has
caveats, particularly the small sample size and differences in software
version, receptor preparation, and sampling parameters between our
protocol and the one the 58\% figure was measured under. The 58\% value
reflects Vina performance in a specific randomized-ligand redocking
experiment and should not be construed as a general ``Vina success rate,''
as other CASF-2016 evaluations using different protocols report
substantially higher top-pose success rates (Supplementary Section~S1).

The CAK series (1/14 passed, 7.1\%; exact 95\%~CI 0.2--33.9\%) showed fewer
successful redockings than the CASF sanity-check subset (8/20, 40\%;
Table~3). Fisher's exact test gave $p=0.037$ for the prespecified one-sided
comparison (hypothesis: CAK would succeed less often than the CASF
control) and $p=0.050$ for the two-sided comparison. The CASF sanity check
did not reveal an obvious gross failure of the docking implementation,
whereas 13 of 14 structures in the sampled CAK series \emph{failed} the
prespecified criterion.

\begin{table}[t]
\centering
\caption{Summary of redocking outcomes}
\label{tab:casf}
\st
\begin{tabular}{lccc}
\toprule
Set & $n$ & Pass (rate) & 95\%~CI \\
\midrule
CASF-2016 subset (pipeline sanity check) & 20 & 8 (40\%)   & 19.1--63.9\% \\
CAK series (single-target)             & 14 & 1 (7.1\%)  & 0.2--33.9\% \\
ATR pair (9L40/9L4B)                   & 2  & 1 (50\%)   & 1.3--98.7\% \\
\bottomrule
\end{tabular}
\end{table}
\FloatBarrier

\subsection{Matched flexible-sidechain analysis}

Allowing flexibility in the five active-site residues nearest the native
berzosertib pose did not recover it: mean symmetry-corrected RMSD across
three seeds was 3.00$\pm$0.92~\AA{} (individual seeds: 4.29, 2.23,
2.49~\AA{}), and no seed met the 2.0~\AA{} threshold. Naively compared to
the original fully rigid result (2.53~\AA{}, 22~\AA{} box, Open
Babel-prepared ligand), it appears that allowing flexibility slightly
worsens the pose -- but this comparison is confounded, since the flexible
run also changes the box size (22~\AA{}$\to$30~\AA{}) and the receptor
preparation (Meeko vs.\ Open Babel) at the same time as flexibility. The
matched control isolates flexibility alone: fully rigid redocking with the
identical enlarged 30~\AA{} box and the identical Meeko-prepared receptor
and ligand, but no residues declared flexible, gives mean
symmetry-corrected RMSD~4.50$\pm$0.05~\AA{} -- worse than both the
original rigid result and the flexible result. Against this correctly
matched baseline, the improvement from 4.50~\AA{} to 3.00~\AA{} when the
five residues are allowed to move is real and in the expected direction,
though still insufficient to reach the 2.0~\AA{} threshold. The larger box
size and the Meeko receptor preparation, not the choice to keep side
chains rigid, account for most of the difference between the original
2.53~\AA{} result and either Meeko-based result. The corrected conclusion
is therefore narrower than either the original claim or its opposite:
flexibility genuinely helps, but not enough on its own to rescue the pose,
so a less-supported deposited pose or a conformational change beyond
these five side chains remains at least as plausible an explanation
(Discussion).

\subsection{Error-propagation comparison}

Docking the 40-compound pilot set against the excluded 9L40 instead of
validated 9L4B produced low rank concordance: Spearman $\rho=0.21$
(bootstrap 95\%~CI $-0.19$ to $0.52$, crossing zero), Kendall~$\tau=0.09$,
mean absolute rank displacement~11.7 positions (max~29), and only 5 of 10
compounds shared between the two top-10 lists (Table~4). The five shared
top-10 compounds were: camonsertib, ceralasertib, deriv\_002, deriv\_003,
and deriv\_019 -- the same five compounds that occupied the top five ranks
in the 9L4B ranking. The top of the ranking is therefore stable across
receptors: a campaign selecting only its top five candidates for synthesis
would reach the same set either way. The churn is concentrated in the
mid-ranks (roughly 6--40), where score gaps are well below Vina's commonly
cited absolute error against experiment (2--3~kcal/mol); a bootstrap
$\rho$ CI that crosses zero describes substantial mid-rank reshuffling,
not a scrambled ranking overall. A reasonable interpretation is that
receptor choice alters the rankings of the middle portion of this pilot
screen but does not affect the highest ranks: a selection from the top
five (or top ten) is relatively insensitive to receptor choice, while a
selection from the top twenty (or top thirty) would draw on a materially
different set of compounds. Again, this is a demonstration of consequence
within this pilot set, not a general enrichment claim, which would require
evaluation against known actives and decoys~\cite{huang2006}.

\begin{table}[t]
\centering
\caption{Largest rank shifts between 9L40 and 9L4B scoring}
\label{tab:errorprop}
\st
\begin{tabular}{lccc}
\toprule
Compound & Rank (9L40) & Rank (9L4B) & Shift \\
\midrule
deriv\_017 & 5  & 34 & 29 \\
deriv\_010 & 10 & 39 & 29 \\
deriv\_022 & 32 & 7  & 25 \\
deriv\_033 & 34 & 9  & 25 \\
\bottomrule
\end{tabular}
\end{table}
\FloatBarrier

\subsection{Implications for practice}

Cognate redocking is a structure-specific quality-control test that can
catch docking errors not apparent from nominal resolution or the evaluated
global map/model metrics. 9L40 and 9L4B were solved in the same study, yet
only the lower-resolution one passed -- a flip that no proxy metric
detected. This observation was reproduced in a separate CAK structure
series, though the CAK structures come from a single target system and so
do not establish generality independent of target. The 13/14 observed
failures in the sampled CAK series should be regarded as descriptive
evidence within that series, not as a kinase-wide failure frequency.

The small size of this survey means that any association seen between the
available global map/model metrics and redocking outcome should not be
over-interpreted. Only two structures with available Q-score measurements
passed, so the comparison of the Q-score group is descriptive rather than
definitive and should not be generalized beyond this sample. The shortcut
to skip the redocking-validation step this paper argues against is not a
neutral one either: the error-propagation comparison above shows that
scoring candidates against an unvalidated receptor does reorder a screen's
mid-ranks even where it preserves its very top, so the real-world cost of
skipping validation depends on how far down a ranked list a downstream
campaign picks candidates. While these results provide evidence for using
per-structure cognate redocking as a quality-control step, the systematic
multi-target evaluation needed to decide on its general utility has yet to
be done.

\subsection{Limitations}

This is a small pilot ($n=17$). The CAK set is a single-target convenience
sample, not a systematic survey of recent kinase depositions. The global
Q-score analysis is underpowered (2 passing structures with reported
Q-scores). Local density-fit metrics were assessed for 9L40 and 9L4B but a
full per-residue active-site analysis across all 17 structures was not
performed, as such an analysis would require uniform local validation-report
data that were not consistently available for the CAK series. No molecular
dynamics, MM-GBSA, or cross-kinase selectivity panel was run. The BRICS
derivative library is small (5 seeds, 35 derivatives) and biased toward
smaller, rule-compliant molecules; enrichment-based screening metrics cannot
be meaningfully computed at this scale. Most importantly, nothing in this
study establishes functional modulation---only predicted binding and ranking.

% ======================================================================
\section{Conclusions}

Cognate redocking is a structure-specific quality-control test that can
reveal failures not obvious from nominal resolution or the evaluated
global map/model metrics. Applying this test to two ATR structures from
the same study showed success on the lower-resolution structure and
failure on the higher-resolution one, an inversion not explained by
allosteric-site proximity. A single-target CAK structure set showed the
cognate-redocking criterion failing in 13 of 14 instances. Its value and
utility for prospective enrichment nonetheless need further assessment
across a wider, multi-target benchmark. Omitting this step is not
innocuous, as shown by this pilot's own error-propagation comparison,
which indicates that scoring candidates against an unvalidated receptor
reorders a screen's mid-ranks, with the degree of impact depending on how
far down a subsequent campaign's candidate selection goes -- even in this
pilot, where the very top ranks happened to be preserved. The next
logical step is to test whether this proportion of failures holds for
recently published kinase structures more broadly, through a systematic
study of many independent targets. The code, data, and structures behind
this work are all available online for independent verification.

% ======================================================================
\section*{Acknowledgements}
[Placeholder]

\section*{Author contributions}
CRediT roles (E.Y.: Ethan Ye; A.P.: Aryan Padarthi; S.K.: Sungju Kim;
A.E.C.: Alexandre Elie-Dit-Cosaque; E.J.: Eunho Jo). Draft assignment,
to be confirmed by all authors prior to submission: Conceptualization --
E.Y.; Methodology -- E.Y.; Software -- E.Y.; Validation -- E.Y., A.P.,
S.K., A.E.C., E.J.; Formal analysis -- E.Y.; Investigation -- E.Y.,
A.P., S.K., A.E.C., E.J.; Data curation -- E.Y.; Writing -- original
draft -- E.Y.; Writing -- review \& editing -- E.Y., A.P., S.K.,
A.E.C., E.J.

\section*{Data availability}
All code, seed and derivative compound tables, redocking validation results,
individual replicate RMSDs, prepared PDBQT files, and full residue-deletion
lists are available at \url{https://github.com/ethanyyy123/research-molecule-editing}
and in the underlying analysis pipeline repository at
\url{https://github.com/The-MathKing/CogRed}, which also tracks every
review round and correction described in this manuscript. Receptor
structures are publicly available from RCSB PDB (9L40, 9L4B, 2YM8, and the
CDK7/cyclin~H/MAT1 and CASF-2016 structures cited in Tables~2--3) and were
retrieved via the AWS Open Data mirror (\url{s3://pdbsnapshots}).

\section*{Declarations}

\textbf{Conflict of interest:} The authors declare no competing interests.

\textbf{Ethical approval:} Not applicable. This study is a purely
computational analysis of publicly deposited structures and did not involve
human participants, human data, or animal subjects.

\textbf{AI assistance:} During the preparation of this work, the authors
used an AI assistant (Claude, Anthropic) to help develop and debug the
computational pipeline's code, run and log individual analysis steps, and
draft and revise sections of this manuscript, and to perform iterative,
automated critique of intermediate drafts against the underlying code and
data -- not a substitute for independent human peer review. That process
itself caught errors, including numbers and described procedures in an
earlier draft that were not actually supported by any script or data file
in the repository above; those have been corrected or removed. All
AI-assisted code, analyses, and text were reviewed by the authors, who take
full responsibility for the content of the published article.

\bibliographystyle{unsrtnat}
\bibliography{refs}

\end{document}